%% =============================================================================
%%  Apresentação — como usar este caderno
%% =============================================================================
\chapter{Apresentação}

Este caderno reúne um estudo pessoal de \rb{Lc 6:20–26}, as quatro felicidades
e os quatro \enquote{ais} que abrem o discurso de Jesus em Lucas~6. O
objetivo é duplo: entender o texto com a maior fidelidade possível e, ao mesmo
tempo, fixar um \emph{método} que possa ser repetido nos próximos cadernos da
série.

\begin{nota}[Sobre o título]
O discurso é conhecido em Lucas como \emph{Sermão da Planície}, e não
\enquote{Sermão do Monte}. Jesus sobe ao monte para orar e escolher os Doze
(\rb{6:12–16}), mas desce com eles e fala \enquote{num lugar plano}
(\gr{ἐπὶ τόπου πεδινοῦ}, \rb{6:17}). O paralelo em Mateus (\rb{Mt 5–7}) é
que se passa no monte. A diferença de cenário é um dado a observar, não um
problema a resolver.
\end{nota}

\section{O roteiro em sete etapas}

O caminho parte do mais próximo do leitor (as traduções que ele já usa) e vai
se aproximando do texto em sua língua, forma e contexto, para só depois ouvir
outros intérpretes e formular uma síntese. Ao final, volta-se ao texto: a
interpretação é uma \emph{espiral}, não uma linha reta \parencite{osborne2006}.

\begin{center}
\begin{tikzpicture}[
    etapa/.style={circle, minimum size=11mm, inner sep=0pt, font=\EBftitulo\Large,
      text=white},
    rot/.style={font=\sffamily\scriptsize, text width=18mm, align=center, text=tinta},
    grupo/.style={font=\EBfversal\footnotesize, text=destaque2!80!black}]
  \foreach \i/\t/\c in {1/Traduções/destaque, 2/Texto original/destaque,
      3/Estrutura/destaque, 4/Contexto/grego!70!destaque2, 5/Intertexto/grego!70!destaque2,
      6/Diálogo/grego!70!destaque2, 7/Síntese/rubrica}{
    \node[etapa, fill=\c] (e\i) at ({(\i-1)*2.1}, 0) {\i};
    \node[rot, below=2pt of e\i] {\t};
  }
  \foreach \i [evaluate=\i as \j using int(\i+1)] in {1,...,6}
    \draw[-{Stealth[length=5pt]}, destaque2, thick] (e\i) -- (e\j);
  \draw[decorate, decoration={brace, amplitude=5pt}, destaque2]
    ([yshift=9mm]e1.west) -- node[grupo, above=6pt]{Observação} ([yshift=9mm]e3.east);
  \draw[decorate, decoration={brace, amplitude=5pt}, destaque2]
    ([yshift=9mm]e4.west) -- node[grupo, above=6pt]{Interpretação} ([yshift=9mm]e6.east);
  \draw[decorate, decoration={brace, amplitude=5pt}, destaque2]
    ([yshift=9mm]e7.west) -- node[grupo, above=6pt]{Apropriação} ([yshift=9mm]e7.east);
  \draw[-{Stealth[length=5pt]}, rubrica, thick, dashed, rounded corners=8pt]
    (e7.south) -- ++(0,-12mm) -| node[pos=.25, below, font=\sffamily\scriptsize, text=rubrica]
    {espiral hermenêutica: voltar ao texto com novas perguntas} (e1.south);
\end{tikzpicture}
\end{center}

Cada etapa abre com uma caixa de \emph{objetivo}, traz perguntas numeradas e
fecha com o que foi aprendido. As etapas retomam o plano original deste
estudo (seis etapas), desdobrando-o para dar lugar próprio ao texto original,
à estrutura literária e ao uso do Antigo Testamento:

\begin{ebtabela}{@{}>{\raggedright}p{.3\linewidth} >{\raggedright\arraybackslash}X@{}}
\EBcab{Plano original} & \EBcab{Neste caderno}\\\midrule
1. Tradução & Etapa 1 — Traduções (comparação crítica)\\
2. Léxico e semântica & Etapa 2 — Texto original (crítica textual, léxico, gramática)\\
3. Contexto imediato & Etapa 3 — Estrutura \emph{e} Etapa 4 — Contexto\\
4. Referências cruzadas & Etapa 5 — Intertexto (AT hebraico, aramaico, LXX, Peshitta)\\
5. Comentários & Etapa 6 — Diálogo com a tradição\\
6. Homilético e doutrinário & Etapa 7 — Síntese e apropriação\\
\end{ebtabela}

\section{Princípios de trabalho}

\paragraph{Observar antes de interpretar.} As três primeiras etapas só
registram o que está no texto: palavras, formas, repetições, contrastes.
Perguntas de sentido são anotadas, mas ficam em aberto
\parencite{duvall2012, fee2002}.

\paragraph{Comentários por último.} Ler comentários cedo demais faz o estudo
repetir conclusões alheias. Eles entram na Etapa~6 como interlocutores, para
confirmar, corrigir ou ampliar o que foi observado.

\paragraph{Palavras têm sentido no contexto.} O interlinear é uma ferramenta de
observação, não uma tradução. Ele tende a sugerir uma leitura \enquote{palavra
por palavra} e a esconder a sintaxe; por isso é sempre acompanhado da análise
da frase inteira.

\begin{alerta}[Falácias a evitar]
\emph{Falácia da raiz}: supor que a etimologia define o sentido atual de uma
palavra. \emph{Transferência ilegítima de totalidade}: carregar para uma
ocorrência todos os sentidos que a palavra tem no dicionário.
\emph{Anacronismo}: ler um termo grego pelo sentido de uma palavra portuguesa
derivada dele. Ver \textcite{carson1996} e \textcite{silva1994}.
\end{alerta}

\section{Legenda}

\noindent
\begin{tabularx}{\linewidth}{@{}l X@{}}
\gr{Grego} & texto grego (Gentium Book Plus, azul)\\
\hb{עִבְרִית} & hebraico (Noto Serif Hebrew, âmbar)\\
\arc{אֲרָמִית} \ \sy{ܣܘܪܝܝܐ} & aramaico e siríaco (verde-oliva)\\
\tl{translit.} & transliteração no padrão SBL\\
\vs{20}\hspace{-.2em} & número de versículo (vermelho-rubrica, como nos manuscritos)\\
{\color{rubrica}°} & palavra com variante textual registrada\\
{\color{suave}[\,]} & palavra implícita, acrescentada na glosa\\
\end{tabularx}

\begin{objetivo}Diz o que se quer alcançar na etapa.\end{objetivo}
\begin{pergunta}Pergunta de estudo, numerada por etapa.\end{pergunta}
\begin{nota}Observação, dado ou esclarecimento.\end{nota}
\begin{variante}Problema de crítica textual.\end{variante}
\begin{aplicacao}Ponte entre o sentido do texto e a vida hoje.\end{aplicacao}


